\section{Metodología}


El procedimiento experimental se estructuró en dos etapas principales orientadas a la caracterización magnética y al análisis de pérdidas en dos configuraciones de núcleo ferromagnético bajo las mismas condiciones de excitación.

\subsection{Núcleo Macizo}
\begin{enumerate}
    \item Se realizó la interconexión de los equipos e instrumentos de medición siguiendo el esquema eléctrico del circuito integrador, empleando para ello un capacitor de $1\,\mu\text{F}$ y una resistencia de $150\,\text{k}\Omega$. Las bobinas disponibles en el laboratorio se acoplaron eléctricamente en configuración serie.
    
    \item Una vez energizado el sistema, se registraron y analizaron mediante el osciloscopio digital las formas de onda correspondientes a la tensión de alimentación y a la corriente.
    
    \item Con ayuda del wattímetro analógico se determinó cuantitativamente la potencia activa total consumida por el conjunto de la bobina y el núcleo ferromagnético macizo.
    
    \item Se procedió a trazar el ciclo de histéresis calibrando los canales del osciloscopio en modo $X-Y$ para un rango de tensiones de alimentación comprendido entre $40\,\text{V}$ y $110\,\text{V}$.
    
    \item  A partir de las dimensiones geométricas del núcleo sólido y el área calculada del ciclo, se determinaron matemáticamente las escalas de campo, las pérdidas totales del núcleo, el componente de pérdidas por histéresis y, por diferencia, las pérdidas debidas a corrientes parásitas. Asimismo, se dedujo analíticamente el valor de la corriente de magnetización para contrastarlo con el valor medido en el instrumento.
\end{enumerate}

\subsection{Núcleo Laminado}
\begin{enumerate}
    \item  Se sustituyó el núcleo macizo por el laminado, manteniendo idénticos los valores del circuito integrador $RC$ ($1\,\mu\text{F}$ y $150\,\text{k}\Omega$) y la conexión en serie de las bobinas de excitación de $4\,\text{A}$.
    
    \item Se observaron en el osciloscopio las variaciones temporales de la tensión de alimentación, la corriente de línea y la señal integrada de tensión.
    
    \item Se registró la lectura del wattímetro para evaluar la nueva potencia activa total absorbida por el núcleo laminado bajo el mismo régimen de excitación.
    
    \item  Se obtuvo el correspondiente lazo de histéresis variando el voltaje de entrada en un rango extendido de $40\,\text{V}$ a $120\,\text{V}$.
    
    \item Se recalcularon las escalas de las variables del lazo magnético, las pérdidas por histéresis y las pérdidas por corrientes parásitas. Para la determinación precisa de la sección transversal efectiva del circuito magnético.
\end{enumerate}